% GENERATED FILE. Edit canonical lessons, then run npm run generate:books.
% canonical-source-hash: fnv1a64:06f2ba6dd8c3c13b
% canonical-lessons: MR-C414-atithi, MR-C414-adaratithya, MR-C414-svagatkaksa, MR-C414-vicarpus-karne, MR-C414-phirakne

\chapter{Guests and Hosts}
\label{ch:mr-guests-and-hosts}
\hlchaptermodality{414}

\begin{chapteropening}
\textbf{By the end of this chapter:} \emph{I can talk about a guest, hospitality, a reception desk, and asking after people.}

Complete the last lesson of chapter 414: I can talk about a guest, hospitality, a reception desk, and asking after people.
\end{chapteropening}

\section[\texorpdfstring{\mr{अतिथी} (atithī)}{अतिथी (atithī)}]{\mr{अतिथी} (atithī) — an honoured guest}

\label{lesson:MR-C414-atithi}



\begin{quote}
Before the new one: say the Marathi for an overnight stay, then the Marathi for an eating house.
\end{quote}

\subsection*{You'll want to know: \mr{अतिथी}}
\textbf{\mr{अतिथी}} — \emph{atithī} — \textquotedblleft{}an honoured guest\textquotedblright{}.

\begin{quote}
\mr{समारंभाला} \mr{मोठे} \mr{अतिथी} \mr{आले} \mr{होते}.
\end{quote}

\begin{quote}
\emph{samārambhālā moṭhe atithī āle hote.}
\end{quote}

\begin{quote}
Distinguished guests had come to the ceremony.
\end{quote}

Masculine.

More formal than \mr{पाहुणा}; the plural is the same: \mr{एक} \mr{अतिथी}, \mr{अनेक} \mr{अतिथी}.

\begin{cousinweb}
\textbf{\mr{अतिथी}} (\emph{atithī}) is a tatsama, taken from Sanskrit \emph{atithi}, \textquotedblleft{}a guest\textquotedblright{} (Monier-Williams).
\end{cousinweb}

\begin{grammarlens}[title={What you've built}]
One more word for getting around.
\end{grammarlens}

\subsection*{Your turn}
\begin{itemize}
\raggedright
  \item \emph{Say it:} \emph{atithī}
  \item \emph{Say it:} \emph{atithī}, once more
  \item \emph{Say it:} \emph{samārambhālā moṭhe atithī āle hote.} — Distinguished guests had come to the ceremony.
  \item \emph{Recall:} say the Marathi for an overnight stay, then the Marathi for an eating house, then say \emph{atithī} again
\end{itemize}

\begin{tcolorbox}[breakable,colback=teal!4,colframe=teal!35!black,title={Before you move on}]
What does \mr{अतिथी} mean? (\textquotedblleft{}An honoured guest\textquotedblright{}.) Say it once more.
\end{tcolorbox}
\section[\texorpdfstring{\mr{आदरातिथ्य} (ādarātithya)}{आदरातिथ्य (ādarātithya)}]{\mr{आदरातिथ्य} (ādarātithya) — hospitality}

\label{lesson:MR-C414-adaratithya}



\begin{quote}
Before the new one: say the Marathi for an eating house, then the Marathi for an honoured guest.
\end{quote}

\subsection*{You'll want to know: \mr{आदरातिथ्य}}
\textbf{\mr{आदरातिथ्य}} — \emph{ādarātithya} — \textquotedblleft{}hospitality\textquotedblright{}.

\begin{quote}
\mr{त्यांच्या} \mr{आदरातिथ्याने} \mr{आम्ही} \mr{खूश} \mr{झालो}.
\end{quote}

\begin{quote}
\emph{tyāñcyā ādarātithyāne āmhī khūś jhālo.}
\end{quote}

\begin{quote}
We were delighted by their hospitality.
\end{quote}

Neuter.

\mr{आदर}, respect, joined to \mr{आतिथ्य}, guest care.

\begin{cousinweb}
\textbf{\mr{आदरातिथ्य}} (\emph{ādarātithya}) is built from \textbf{\mr{आदर}}, \textquotedblleft{}respect\textquotedblright{}, and Sanskrit \emph{ātithya}, \textquotedblleft{}proper for a guest, hospitality\textquotedblright{} (Monier-Williams).
\end{cousinweb}

\begin{grammarlens}[title={What you've built}]
One more word for getting around.
\end{grammarlens}

\subsection*{Your turn}
\begin{itemize}
\raggedright
  \item \emph{Say it:} \emph{ādarātithya}
  \item \emph{Say it:} \emph{ādarātithya}, once more
  \item \emph{Say it:} \emph{tyāñcyā ādarātithyāne āmhī khūś jhālo.} — We were delighted by their hospitality.
  \item \emph{Recall:} say the Marathi for an eating house, then the Marathi for an honoured guest, then say \emph{ādarātithya} again
\end{itemize}

\begin{tcolorbox}[breakable,colback=teal!4,colframe=teal!35!black,title={Before you move on}]
What does \mr{आदरातिथ्य} mean? (\textquotedblleft{}Hospitality\textquotedblright{}.) Say it once more.
\end{tcolorbox}
\section[\texorpdfstring{\mr{स्वागतकक्ष} (svāgatkakṣa)}{स्वागतकक्ष (svāgatkakṣa)}]{\mr{स्वागतकक्ष} (svāgatkakṣa) — a reception desk}

\label{lesson:MR-C414-svagatkaksa}



\begin{quote}
Before the new one: say the Marathi for an honoured guest, then the Marathi for hospitality.
\end{quote}

\subsection*{You'll want to know: \mr{स्वागतकक्ष}}
\textbf{\mr{स्वागतकक्ष}} — \emph{svāgatkakṣa} — \textquotedblleft{}a reception desk\textquotedblright{}.

\begin{quote}
\mr{स्वागतकक्षात} \mr{आम्हाला} \mr{खोलीची} \mr{किल्ली} \mr{मिळाली}.
\end{quote}

\begin{quote}
\emph{svāgatkakṣāt āmhālā kholīcī killī miḷālī.}
\end{quote}

\begin{quote}
We got our room key at the reception desk.
\end{quote}

Masculine.

The front desk of a hotel or an office; \mr{कक्ष} is a room.

\begin{cousinweb}
\textbf{\mr{स्वागतकक्ष}} (\emph{svāgatkakṣa}) is a modern word built from Sanskrit \emph{svāgata}, \textquotedblleft{}welcome\textquotedblright{}, and \emph{kakṣa}, \textquotedblleft{}an inner room\textquotedblright{} (Monier-Williams).
\end{cousinweb}

\begin{grammarlens}[title={What you've built}]
One more word for getting around.
\end{grammarlens}

\subsection*{Your turn}
\begin{itemize}
\raggedright
  \item \emph{Say it:} \emph{svāgatkakṣa}
  \item \emph{Say it:} \emph{svāgatkakṣa}, once more
  \item \emph{Say it:} \emph{svāgatkakṣāt āmhālā kholīcī killī miḷālī.} — We got our room key at the reception desk.
  \item \emph{Recall:} say the Marathi for an honoured guest, then the Marathi for hospitality, then say \emph{svāgatkakṣa} again
\end{itemize}

\begin{tcolorbox}[breakable,colback=teal!4,colframe=teal!35!black,title={Before you move on}]
What does \mr{स्वागतकक्ष} mean? (\textquotedblleft{}A reception desk\textquotedblright{}.) Say it once more.
\end{tcolorbox}
\section[\texorpdfstring{\mr{विचारपूस} \mr{करणे} (vicārpūs karṇe)}{विचारपूस करणे (vicārpūs karṇe)}]{\mr{विचारपूस} \mr{करणे} (vicārpūs karṇe) — to ask after}

\label{lesson:MR-C414-vicarpus-karne}



\begin{quote}
Before the new one: say the Marathi for hospitality, then the Marathi for a reception desk.
\end{quote}

\subsection*{You'll want to know: \mr{विचारपूस} \mr{करणे}}
\textbf{\mr{विचारपूस} \mr{करणे}} — \emph{vicārpūs karṇe} — \textquotedblleft{}to ask after\textquotedblright{}.

\begin{quote}
\mr{आजीने} \mr{फोनवर} \mr{सगळ्यांची} \mr{विचारपूस} \mr{केली}.
\end{quote}

\begin{quote}
\emph{ājīne phonvar sagaḷyāñcī vicārpūs kelī.}
\end{quote}

\begin{quote}
Grandmother asked after everyone on the phone.
\end{quote}

Past, with -\mr{ने} (he …): \mr{त्याने} \mr{विचारपूस} \mr{केली}, \emph{tyāne vicārpūs kelī}.

The person asked about takes -\mr{ची}: \mr{माझी} \mr{विचारपूस} \mr{केली}, asked after me.

\begin{cousinweb}
\textbf{\mr{विचारपूस} \mr{करणे}} (\emph{vicārpūs karṇe}) is built from \textbf{\mr{विचार}}, \textquotedblleft{}a thought, a question\textquotedblright{}, and \emph{pūs}, \textquotedblleft{}asking\textquotedblright{}, from the family of Sanskrit \emph{pṛcchati}, \textquotedblleft{}asks\textquotedblright{}, where CDIAL lists Marathi \emph{pusṇẽ}, \textquotedblleft{}to ask\textquotedblright{} (CDIAL 8352); \textbf{\mr{करणे}} makes it a verb.
\end{cousinweb}

\begin{grammarlens}[title={What you've built}]
One more word for getting around.
\end{grammarlens}

\subsection*{Your turn}
\begin{itemize}
\raggedright
  \item \emph{Say it:} \emph{vicārpūs karṇe}
  \item \emph{Say it:} \emph{vicārpūs karṇe}, once more
  \item \emph{Say it:} \emph{ājīne phonvar sagaḷyāñcī vicārpūs kelī.} — Grandmother asked after everyone on the phone.
  \item \emph{Recall:} say the Marathi for hospitality, then the Marathi for a reception desk, then say \emph{vicārpūs karṇe} again
\end{itemize}

\begin{tcolorbox}[breakable,colback=teal!4,colframe=teal!35!black,title={Before you move on}]
What does \mr{विचारपूस} \mr{करणे} mean? (\textquotedblleft{}To ask after\textquotedblright{}.) Say it once more.
\end{tcolorbox}
\section[\texorpdfstring{\mr{फिरकणे} (phirakṇe)}{फिरकणे (phirakṇe)}]{\mr{फिरकणे} (phirakṇe) — to drop by}

\label{lesson:MR-C414-phirakne}



\begin{quote}
Before the new one: say the Marathi for a reception desk, then the Marathi for to ask after.
\end{quote}

\subsection*{You'll want to know: \mr{फिरकणे}}
\textbf{\mr{फिरकणे}} — \emph{phirakṇe} — \textquotedblleft{}to drop by\textquotedblright{}.

\begin{quote}
\mr{सुट्टीत} \mr{तो} \mr{आमच्याकडे} \mr{फिरकलाच} \mr{नाही}.
\end{quote}

\begin{quote}
\emph{suṭṭīt to āmcyākaḍe phiraklāc nāhī.}
\end{quote}

\begin{quote}
In the holidays he did not even drop by our place.
\end{quote}

Past (he …): \mr{तो} \mr{फिरकला}, \emph{to phiraklā}.

Mostly with a negative: \mr{इकडे} \mr{फिरकला} \mr{नाही}, he never came this way.

\begin{cousinweb}
\textbf{\mr{फिरकणे}} (\emph{phirakṇe}) is a Marathi verb that looks built on \emph{phir-}, \textquotedblleft{}to turn\textquotedblright{} (CDIAL 9078); how it came to mean dropping in is not settled.
\end{cousinweb}

\begin{grammarlens}[title={What you've built}]
One more word for getting around.
\end{grammarlens}

\subsection*{Your turn}
\begin{itemize}
\raggedright
  \item \emph{Say it:} \emph{phirakṇe}
  \item \emph{Say it:} \emph{phirakṇe}, once more
  \item \emph{Say it:} \emph{suṭṭīt to āmcyākaḍe phiraklāc nāhī.} — In the holidays he did not even drop by our place.
  \item \emph{Recall:} say the Marathi for a reception desk, then the Marathi for to ask after, then say \emph{phirakṇe} again
\end{itemize}

\begin{tcolorbox}[breakable,colback=teal!4,colframe=teal!35!black,title={Before you move on}]
What does \mr{फिरकणे} mean? (\textquotedblleft{}To drop by\textquotedblright{}.) Say it once more.
\end{tcolorbox}
